\documentclass[svgnames]{article}     % use "amsart" instead of "article" for AMSLaTeX format
%\geometry{landscape}                 % Activate for rotated page geometry

%\usepackage[parfill]{parskip}        % Activate to begin paragraphs with an empty line rather than an indent

\usepackage{graphicx}                 % Use pdf, png, jpg, or eps§ with pdflatex; use eps in DVI mode

%maths                                % TeX will automatically convert eps --> pdf in pdflatex
\usepackage{amssymb}
\usepackage{amsmath}
\usepackage{esint}
\usepackage{geometry}

% Inverting Color of PDF
%\usepackage{xcolor}
%\pagecolor[rgb]{0.19,0.19,0.19}
%\color[rgb]{0.77,0.77,0.77}

%noindent
%\setlength\parindent{0pt}

%pgfplots
\usepackage{pgfplots}

%images
%\graphicspath{{ }}                   % Activate to set a image directory

%tikz
\usepackage{pgfplots}
\pgfplotsset{compat=1.15}
\usepackage{comment}
\usetikzlibrary{arrows}
\usepackage[most]{tcolorbox}

%Figures
\usepackage{float}
\usepackage{caption}
\usepackage{lipsum}


\title{Nondegenerate Perturbation Theory}
\author{Deval Deliwala}
%\date{}                              % Activate to display a given date or no date

\begin{document}
\maketitle
%\section{}
%\subsection{}
%\tableofcontents                     % Activate to display a table of contents


\section{General Formulation} 

Provided a set of orthonormal solutions to the (time-independent) Schr\"odinger
equation for some potential, (say, a 1D infinite square well):

\[
H^0 \psi_n^0 = E_n^0 \psi_n^0 \quad \Leftrightarrow \quad \langle \psi_n^0
| \psi_m^0\rangle = \delta_{nm}
\] \vspace{3px}

with corresponding energy eigenvalues $E_n^0$. Now we \textit{perturb} the
potential slightly, say, by putting a small bump in the bottom of the well.
This new potential is much more complicated and solving it would be difficult
to say the least. \textbf{Perturbation Theory} is a systematic procedure for
obtaining \textit{approximate solutions}  to the perturbed problem, by building
on the known exact solutions to the unperturbed case. 

To begin with, we write the new Hamiltonian as the sum of two terms: 

\[
H = H^0 + \lambda H' 
\] \vspace{3px}

where $H'$ is the perturbation. For the moment we will take $\lambda$ to be
a small number and $H$ is the true Hamiltonian. Next, we write $\psi_n$ and
$E_n$ as a power series in $\lambda$ : 

\begin{align}
  \psi_n &= \psi_n^0 + \lambda \psi_n^1 + \lambda^2 \psi_n^2 + \cdots, \label{1} \\
  E_n &= E_n^0 + \lambda E_n^1 + \lambda^2 E_n^2 + \cdots,  \label{2}
\end{align}

Here, $E_n^1$ is the \textbf{first-order correction} to the $n $th
eigenvalue, and $\psi_n^1$ is the first-order correction to the $n$th
eigenfunctions. Plugging Eqs. (\ref{1}, \ref{2}) into the Hamiltonian
eigenvalue problem, we get 

\begin{align*} \label{}
  H\psi_n &= E_n \psi_n, \\
  (H^0+ \lambda H') \left[ \psi_n^0 + \lambda \psi_n^1 + \lambda^2 \psi_n^2
+ \cdots \right] &= \\ &\left(E_n^0 + \lambda E_n^1 + \lambda^2 E_n^2 + \cdots
\right) \left[\psi_n^0 + \lambda \psi_n^1 + \lambda^2\psi_n^2 + \cdots \right]
\end{align*}\vspace{3px}

Collecting like powers of $\lambda$, 

\begin{align*}
  H^0 &\psi_n^0 + \lambda (H^0 \psi_n^1 + H'\psi_n^0) + \lambda^2(H^0
  \lambda_n^2 + H' \psi_n^1) + \cdots \\ &= E_n^0 \psi_n^0 + \lambda
  (E_n^0\psi_n^1 + E_n^1 \psi_n^0) + \lambda^2 (E_n^0 \psi_n^2 + E_n^1 \psi_n^1
  + E_n^2\psi_n^0) + \cdots.
\end{align*}

To the lowest order ($\lambda^0$ ), this yields $H^0\psi_n^0 = E_n^0 \psi_n^0$,
which is nothing new. To first order ($\lambda^1$), 

 \begin{align} \label{first}
  H^0\psi_n^1 + H'\psi_n^0 = E_n^0\psi_n^1 + E_n^2\psi_n^0
\end{align}\vspace{3px}

To second order ($\lambda^2$), 

\begin{align} \label{second}
  H^0\psi_n^2 + H'(\psi_n^1) = E_n^0\psi_n^2 + E_n^1 \psi_n^1 + E_n^2 \psi_n^0
\end{align}

and so on. Now we are done with $\lambda$, it was just a device to keep track
of the different orders. 

\section{First-Order Theory} 

We will first focus on first-order perturbation, defined by Eq. (\ref{first}).
Taking the inner product with $\psi_n^0$, (multiplying by $(\psi_n^0)^*$ and
integrating), 

\[
\langle \psi_n^0 | H^0 \psi_n^1 \rangle + \langle \psi_n^0 | H' \psi_n^0
\rangle = E_n^0 \langle \psi_n^0 | \psi_n^1\rangle + E_n^1 \langle \psi_n^0
| \psi_n^0 \rangle. 
\] \vspace{3px}

But since $H^0$ is hermitian, we get 

\[
\langle \psi_n^0 | H^0 | \psi_n^1\rangle = \langle H^0 \psi_n^0 | \psi_n^1
\rangle = \langle E_n^0\psi_n^0 | \psi_n^1\rangle = E_n^0\langle \psi_n^0
| \psi_n^1 \rangle. 
\] \vspace{3px}

Therefore, cancelling the first term on the right and setting $\langle \psi_n^0
| \psi_n^0 \rangle = 1$, we get 

\vspace{3px}
\begin{tcolorbox}[colback = red!5!white, colframe = red!50!black, title = First
  Order Correction to Energy]

  \begin{align}\label{firstenergy} E_n^1 = \langle \psi_n^0 | H' | \psi_n^0
  \rangle. \end{align}

\end{tcolorbox}
\vspace{3px}

This is the fundamental result of first-order perturbation theory, likely the
most used equation in quantum mechanics. The first-order correction to the
energy is the \textit{expectation value} of the perturbation, in the
\textit{unperturbed} state. 


Eq. (\ref{firstenergy}) is the first-order correction to the \textit{energy}.
To find the first-order correction to the \textit{wave function}, we rewrite
Eq. (\ref{first}): 

\begin{align} \label{b}
  H^0\psi_n^1 + H'\psi_n^0 &= E_n^0\psi_n^1 + E_n^2\psi_n^0 \nonumber \\ 
  \left(H^0 - E_n^0\right) \psi_n^1 &= - \left( H' - E_n^1 \right) \psi_n^0. 
\end{align}

The right side is a known function, so this amounts to an inhomogeneous
differential equation for $\psi_n^1$. Now the unperturbed wave functions
constitute a complete set, so $\psi_n^1$ can be expressed as a linear
combination of them: 

\begin{align} \label{a}
  \psi_n^1 = \sum_{m\neq n}^{} c_m^{(n)} \psi_m^0
\end{align}\vspace{3px}

Putting Eq. (\ref{a}) into Eq. (\ref{b}) and using the fact taht $\psi_m^0$
satisfies the unperturbed Schr\"odinger equation, we have 

\[
\sum_{m\neq n}^{} \left( E_m^{0} - E_n^{0}  \right) c_m^{(n)} \psi_m^{0}
= -\left( H' - E_n^{1} \right) \psi_n^{0}.  
\] \vspace{3px}

Taking the inner product now with $\psi_\ell^{0}$, 

\[
\sum_{m\neq n}^{} \left( E_m^{0} - E_n^{0} \right) c_m^{(n)} \langle \psi_\ell^{0}
| \psi_m^{0} = -\langle \psi_\ell^{0} | H' | \psi_n^{0} \rangle = E_n^1 \langle
\psi_\ell^{0} | \psi_n^{0} \rangle.
\] \vspace{3px}

If $\ell = n$, the left side is 0, and we recover Eq. (\ref{firstenergy}). If
$\ell \neq n$, we get 

\[
\left( E_\ell^{0} - E_n^{0}  \right) c_\ell^{(n)} = -\langle \psi_\ell^{0} | H'
| \psi_n^{0} \rangle,
\] \vspace{3px}

resulting in 

\begin{align} \label{coeff}
  c_m^{(n)} = \frac{\langle \psi_m^{0} | H' | \psi_n^0\rangle}{E_n^{0}
  - E_m^{0}}.
\end{align}\vspace{3px}

Therefore, 

\vspace{3px}
\begin{tcolorbox}[colback = blue!5!white, colframe = blue!50!black, title
  = First Order Correction to Wave Function]

  \begin{align} \label{firstwf}
    \psi_n^{1} = \sum_{m\neq n} \frac{{\langle \psi_m^{0}| H'  |\psi_n^0
    \rangle} }{(E_n^0 - E_m^0)} \psi_m^0.
  \end{align}

\end{tcolorbox}
\vspace{3px}
The denominator is safe (since there is no coefficient with $m = n$
) \textit{as long as the unperturbed energy spectrum is nondegenerate}. If two
different unperturbed states share the same energy, we divided by zero to get
Eq. (\ref{coeff}). In this case we need \textbf{degenerate perturbation
theory}. 

This completes first-order perturbation theory. The first-order correction to
the energy, $E_n^1$, is given by Eq. (\ref{firstenergy}), and the first-order
correction to the wave function, $\psi_n^1$, is given by Eq. (\ref{firstwf}). 





\end{document}

